\section{Conclusion}
\label{sec:conclusion}

MorphoCLIP maps Cell Painting wells and text descriptions of perturbations into a 512-dimensional space with about 14M trainable parameters on top of two frozen backbones. One model handles compounds, CRISPR knockouts and ORF overexpressions. At one seed, perturbation-level Recall@10 is about 39\% in both directions on validation and 37 to 44\% on test, against analytic chance rates of 10 to 12\%. Separate stochastic benchmark runs show numerically comparable compound replicability fractions for MorphoCLIP and the released CellCLIP checkpoint but do not support a ranking. The base model's four-track mean compound target-matching fraction lies within the published CellProfiler range, although one track is above it.

Gene-aware soft labels, the replicate alignment loss and condition-relative plate offsets each change validation retrieval by only a few perturbations, and the test split reorders them. These retrieval differences are unresolved without repeated seeds. The replicate loss does increase deterministic replicability mAP on every track, from pooled 0.298 to 0.343, while the combined run reaches 0.369. This result is partly definitional because the loss directly optimizes replicate proximity. The in-batch plate correction degrades retrieval because the plate mean carries condition information named by the prompt; the condition-relative offset avoids that degradation but shows no consistent retrieval or benchmark-mAP improvement at this seed count.

Gene--compound fraction retrieved is at or near zero for every model: at most three admitted pairs pass the stochastic significance gate in any run. We have no positive cross-modality result.

\subsection{Limitations and Future Work}

Every retrieval number in this paper is from one seed. Repeated seeds are needed for the control, replicate-loss and combined runs. Fraction retrieved on the standard benchmark varies between runs of the same code, so those values are descriptive. We did not rerun CellProfiler, and the CellCLIP row comes from a separate run with a different export path and a batch-correction step ours lacks. Because the replicate loss directly optimizes the replicability task, evidence that the embedding is more informative requires a task outside that objective; target matching does not provide such evidence here.

On cross-modality, one concrete lead is the choice of gene annotation. Training soft labels use each compound's primary target (160 distinct genes) while the benchmark's gene--compound track uses the full target list (758 distinct genes, off-targets included). Aligning the two, or training on the wider list, is a one-line change we have not measured. A second lead is prompt content: whether the SMILES string and gene-function text carry signal through BioClinical ModernBERT, or are inert, can be tested by ablating them out of the templates. Beyond CPJUMP1, the full JUMP corpus of about 140,000 perturbations across many sites and cell lines is the setting in which a text-supervised model would have to prove that its embeddings transfer.

We release the code, configs, split manifests and the evaluation harness with the paper.
